\chapter{System Architecture and Methodology}
\label{ch:methodology}
% Chapter 4: System Architecture and Methodology (15-18 pages)

%
% Section: System Overview
%
\section{System Overview}
\label{sec:methodology:overview}

% Prof. Andreas: "4.1 Sys Overview ← should contain picture (Lane1 / Lane2)"
% System architecture diagram showing the two-lane retrieval + generator pipeline

This section presents an overview of the proposed multimodal RAG system for Leishmaniasis diagnosis. The architecture follows a two-lane design that separates text-based retrieval from vision-language retrieval, allowing each modality to leverage domain-specific encoders optimized for medical content.

\begin{figure}[htbp]
  \centering
  % TODO: Replace with actual architecture diagram
  % Suggested diagram structure:
  % Lane 1: Query Text → E5-large-instruct → Dense Retrieval → RRF Fusion ↘
  %                    → BM25 Sparse Retrieval →                         ↗  → Retrieved Docs
  % Lane 2: Query Image → BiomedCLIP → Cross-modal Retrieval →          ↗
  % Retrieved Docs → Context Builder → LLM (MedGemma/Gemini) → Diagnosis
  \fbox{\parbox{0.9\textwidth}{\centering\vspace{3cm}
    \textbf{[System Architecture Diagram Placeholder]}\\[1em]
    Two-Lane Retrieval Architecture:\\
    Lane 1: Text Retrieval (E5-large-instruct + BM25 → RRF Fusion)\\
    Lane 2: Vision-Language Retrieval (BiomedCLIP)\\
    → Context Integration → Answer Generation (LLM)
  \vspace{3cm}}}
  \caption[Two-Lane RAG Architecture]{Overview of the two-lane multimodal RAG architecture. Lane 1 handles text-based retrieval using E5-large-instruct dense embeddings combined with BM25 sparse retrieval via Reciprocal Rank Fusion (RRF). Lane 2 processes visual queries using BiomedCLIP for cross-modal image-text matching. Retrieved contexts from both lanes are integrated and passed to the generator (MedGemma or Gemini) for diagnosis generation.}
  \label{fig:methodology:architecture}
\end{figure}

\subsection{Design Goals and Requirements}
\label{subsec:methodology:overview:goals}

\subsection{Architecture Philosophy}
\label{subsec:methodology:overview:philosophy}
% Two-lane approach rationale

%
% Section: Retriever Model Specifications
% Prof. Andreas: "retriever model . tex" - describe retriever model clearly
%
\section{Retriever Model Specifications}
\label{sec:methodology:retriever}

% Prof. Andreas: Document model parameters (size, dim, RAG settings) for fair comparison
% "models same parameters / model size / dim / rag"

This section documents the retriever models used in our system with their key specifications for reproducibility and fair comparison.

\begin{table}[htbp]
  \centering
  \caption{Retriever Model Specifications}
  \label{tab:methodology:retriever-specs}
  \begin{tabularx}{\textwidth}{lXX}
    \toprule
    \textbf{Component} & \textbf{Lane 1 (Text)} & \textbf{Lane 2 (Vision-Language)} \\
    \midrule
    Model & E5-large-instruct & BiomedCLIP \\
    Base Architecture & BERT-large & ViT-B/16 + BioBERT \\
    Embedding Dimension & 1024 & 512 \\
    Model Size & 335M parameters & 151M parameters \\
    Pooling Strategy & Mean pooling & [CLS] token \\
    Normalization & L2 normalized & L2 normalized \\
    Index Type & HNSW (Qdrant) & HNSW (Qdrant) \\
    Sparse Component & BM25 (k1=1.5, b=0.75) & N/A \\
    Fusion Method & RRF (k=60) & Direct similarity \\
    \bottomrule
  \end{tabularx}
\end{table}

%
% Section: Lane 1 - Text Retrieval
%
\section{Lane 1: Text Retrieval}
\label{sec:methodology:lane1}

\subsection{E5-large-instruct Encoder}
\label{subsec:methodology:lane1:e5}
% 1024-dimensional embeddings

\subsection{BM25 Sparse Retrieval}
\label{subsec:methodology:lane1:bm25}

\subsection{Hybrid Fusion with RRF}
\label{subsec:methodology:lane1:rrf}
% Reciprocal Rank Fusion

%
% Section: Lane 2 - Vision-Language Retrieval
%
\section{Lane 2: Vision-Language Retrieval}
\label{sec:methodology:lane2}

\subsection{BiomedCLIP Model}
\label{subsec:methodology:lane2:biomedclip}
% 512-dimensional embeddings

\subsection{Cross-modal Similarity Computation}
\label{subsec:methodology:lane2:similarity}

\subsection{Image-Query Matching}
\label{subsec:methodology:lane2:matching}

%
% Section: Index Construction
%
\section{Index Construction}
\label{sec:methodology:indexing}

\subsection{Qdrant Vector Database}
\label{subsec:methodology:indexing:qdrant}

\subsection{Collection Design}
\label{subsec:methodology:indexing:collections}
% Versioning (v143)

\subsection{Stable Hashing for Reproducibility}
\label{subsec:methodology:indexing:hashing}

%
% Section: Retrieval Pipeline
%
\section{Retrieval Pipeline}
\label{sec:methodology:retrieval}

\subsection{Query Processing}
\label{subsec:methodology:retrieval:query}

\subsection{Two-Lane Search Execution}
\label{subsec:methodology:retrieval:execution}

\subsection{Result Aggregation}
\label{subsec:methodology:retrieval:aggregation}

%
% Section: Answer Generation
%
\section{Answer Generation}
\label{sec:methodology:generation}

\subsection{Prompt Engineering}
\label{subsec:methodology:generation:prompts}
% Augmentation-based approach

\subsection{Context Integration}
\label{subsec:methodology:generation:context}

\subsection{Generator Model Options}
\label{subsec:methodology:generation:models}
% Gemini, MedGemma, etc.

%
% Section: Multimodal Integration
%
\section{Multimodal Integration}
\label{sec:methodology:multimodal}

\subsection{Image Embedding Strategy}
\label{subsec:methodology:multimodal:embedding}

\subsection{Text-Image Fusion}
\label{subsec:methodology:multimodal:fusion}

\subsection{Visual Context in Generation}
\label{subsec:methodology:multimodal:visual}
